\chapter{Esercitazione E2: programmazione di rete}
\titolobreve{Esercitazione E2}

Il secondo pacchetto: socket, processi, thread, sincronizzazione, server concorrenti, multiplexing, framing e UDP. L'esercizio $n$ di
questo capitolo è l'\textbf{E2.$n$} delle slide. Tre tipi: \emph{su carta} (leggere codice e dire che cosa fa, o trovare che cosa non
va), \emph{da scrivere} (un programma C completo, con il modo di provarlo), \emph{da catturare} (eseguirlo e guardare il filo). Nelle
slide il codice omette i controlli d'errore, salvo quando l'errore è il punto; il vostro non dovrebbe. Lo stesso vale per i frammenti di
soluzione qui sotto, che mostrano solo la parte che conta.

\info{Ambiente e prove}{
    Va bene qualsiasi Linux, anche una macchina virtuale. Si compila con gli avvisi accesi, e con i thread dove servono:
    \texttt{gcc -Wall -Wextra -pthread -o server server.c}. Per provare un server TCP o UDP il primo client è \texttt{nc};
    \texttt{ss -tan} e \texttt{ss -uan} mostrano chi è legato a che cosa e chi è connesso. Su una sola macchina si cattura sull'interfaccia
    di loopback: i pacchetti tra i due programmi non raggiungono mai un cavo.
}

% ══════════════════════════════════════════════════════════════
\section{Socket e ordine dei byte}
% ══════════════════════════════════════════════════════════════

\begin{esercizio}{Quanti socket}
Un server UDP ha bisogno di un socket, un server TCP di due. Perché? Se il server TCP serve $n$ client su $n$ host diversi
contemporaneamente, quanti socket ha aperti?

\svolgimento
UDP non ha connessioni: un solo socket riceve i datagrammi di tutti i client e risponde all'indirizzo da cui ciascuno arriva. In TCP il
socket in ascolto accetta soltanto le connessioni, e \texttt{accept} crea un socket \textbf{connesso} per ciascun client, identificato
dalla quadrupla degli indirizzi. Con $n$ client: $n + 1$ socket.
\end{esercizio}

\begin{esercizio}{Il client prima del server}
Con i client e server TCP e UDP scritti per le lezioni sui socket: si avvia ogni client prima del suo server. Che cosa succede, e perché
è diverso per i due protocolli? Poi si avviano entrambi, con il client che usa un numero di porta diverso da quello del server.

\svolgimento
TCP: \texttt{connect} fallisce subito con \texttt{ECONNREFUSED}, perché all'arrivo del SYN il kernel del server risponde con un RST.
UDP: \texttt{sendto} riesce (non c'è handshake), il datagramma arriva e viene scartato, e il client resta bloccato in \texttt{recvfrom}
per sempre (l'ICMP \emph{port unreachable} che torna non gli viene riferito, a meno che il socket non sia connesso). Con la porta
sbagliata succede esattamente lo stesso: per il kernel del server su quella porta non c'è nessuno.
\end{esercizio}

\begin{esercizio}{Un bind nel client UDP}
Nel client UDP si aggiunge \texttt{bind} sulla porta 5432 subito dopo \texttt{socket}. Il server deve cambiare? Quali porte usano i due
socket prima e dopo la modifica? Verificarlo con \texttt{ss -uan}.

\svolgimento
Il server non cambia: risponde all'indirizzo e alla porta sorgente che legge da \texttt{recvfrom}. Prima il client riceve una porta
\textbf{effimera} scelta dal kernel al primo \texttt{sendto} (diversa a ogni esecuzione); dopo usa sempre 5432. La porta del server non
cambia. Se si avviano due client insieme, il secondo \texttt{bind} fallisce con \texttt{EADDRINUSE}.
\end{esercizio}

\begin{esercizio}{Ordine dei byte}
Un client su x86 invia un contatore; il server gira su una macchina big-endian e stampa ciò che legge.
\begin{lstlisting}[language={[POSIX]C}]
uint32_t count = 1000;                       /* client: 0x000003E8 */
write(fd, &count, sizeof count);

uint32_t got;                                /* server */
read(fd, &got, sizeof got);                  /* tutti e 4 i byte arrivano */
printf("%u\n", got);
\end{lstlisting}
\begin{enumerate}[label=(\alph*)]
    \item Scrivere i quattro byte nell'ordine in cui viaggiano. Che cosa stampa il server?
    \item Correggere entrambi i lati. La correzione funziona ancora se le due macchine sono entrambe x86?
    \item Quali due byte mette sul filo \texttt{htons(5000)}, e in che ordine?
\end{enumerate}

\svolgimento
\begin{enumerate}[label=(\alph*)]
    \item x86 è little-endian: viaggiano \texttt{E8 03 00 00}. Il server big-endian li legge come \texttt{0xE8030000} e stampa
          3\,892\,510\,720.
    \item Il client invia \texttt{htonl(count)}, il server stampa \texttt{ntohl(got)}. Sì: su x86 le due funzioni scambiano i byte
          entrambe, e sul filo c'è sempre l'ordine di rete (big-endian); su una macchina big-endian non fanno nulla.
    \item $5000 = \texttt{0x1388}$: viaggiano \texttt{13} e poi \texttt{88}.
\end{enumerate}
\end{esercizio}

\begin{esercizio}{Una read non è un messaggio}
Un client chiede a un server un file di 3000 byte e lo stampa.
\begin{lstlisting}[language={[POSIX]C}]
char buf[4096];
write(fd, "GET file\n", 9);
ssize_t n = read(fd, buf, sizeof buf);
printf("%s", buf);
close(fd);
\end{lstlisting}
\begin{enumerate}[label=(\alph*)]
    \item Il file arriva completo quando client e server sono sulla stessa macchina, e a volte tagliato quando non lo sono. Perché?
    \item Indicare un secondo difetto, indipendente, nella riga del \texttt{printf}.
    \item Riscrivere la lettura in modo che funzioni qualunque cosa faccia la rete. Come sa il client che il file è finito?
\end{enumerate}

\svolgimento
\begin{enumerate}[label=(\alph*)]
    \item TCP è un flusso di byte: \texttt{read} restituisce ciò che è arrivato finora. Su loopback i 3000 byte arrivano in un colpo, in
          rete possono arrivare in più segmenti, e la prima \texttt{read} ne restituisce solo una parte.
    \item \texttt{buf} non è terminato da \verb|'\0'|, quindi \texttt{printf} legge oltre i byte ricevuti; e se il file contiene byte nulli,
          \texttt{\%s} si ferma al primo. Inoltre \texttt{n} non viene controllato (potrebbe essere $-1$).
    \item Si legge in un ciclo e si scrive esattamente ciò che si è letto:
\begin{lstlisting}[language={[POSIX]C}]
ssize_t n;
while ((n = read(fd, buf, sizeof buf)) > 0)
    fwrite(buf, 1, (size_t) n, stdout);
if (n < 0) perror("read");
\end{lstlisting}
          Il file è finito quando \texttt{read} restituisce 0: il server ha chiuso la connessione. In alternativa il server invia prima la
          lunghezza, e il client legge esattamente quei byte.
\end{enumerate}
\end{esercizio}

\begin{esercizio}{Trovare gli errori}
Questo server con \texttt{fork} viene da una vecchia dispensa. È sbagliato in molti punti: alcuni impediscono la compilazione, altri si
comportano male a tempo di esecuzione. Trovare almeno otto errori: per ciascuno la riga, che cosa va storto e la correzione.
\begin{lstlisting}[language={[POSIX]C}, style=wnumbers]
#define MAX_DIM 1024
#define MAX_CODA 3

void main() {
    int ds_sock, ds_sock_a;
    struct sockaddr_in server;
    struct sockaddr client;
    char buff[MAX_DIM];

    sock = socket(AF_INET, SOCK_STREAM, 0);
    server.sin_family      = AF_INET;
    server.sin_port        = 1999;
    server.sin_addr.s_addr = INADDR_ANY;
    bind(ds_sock, &server, sizeof(server));
    listen(ds_sock, MAX_CODA);
    length = sizeof(client);
    signal(SIGCHLD, SIG_IGN);
    while (1) {
        while ((ds_sock_a = accept(ds_sock, &client, &length)) == -1);
        if (fork() == 0) {
            do {
                 read(ds_sock, buff, MAX_DIM);
                 printf("client says: %s\n", buff);
            } while (strcmp(buff, "quit") != 0);
            write(ds_sock, "done", 10);
            close(ds_sock_a);
            exit(0);
        } else close(ds_sock_a);
    }
}
\end{lstlisting}

\svolgimento
\begin{itemize}
    \item Riga 4: \texttt{void main()} non è standard: \texttt{int main(void)}, e mancano gli \texttt{\#include}.
    \item Riga 10: \texttt{sock} non è dichiarata: il risultato va in \texttt{ds\_sock}, che altrimenti resta indefinito.
    \item Riga 12: la porta va in ordine di rete, \texttt{htons(1999)}; e la struttura andrebbe azzerata con \texttt{memset} prima.
    \item Riga 13: \texttt{htonl(INADDR\_ANY)} (vale 0, quindi funziona per caso).
    \item Riga 14: serve il cast \texttt{(struct sockaddr *) \&server}; e nessuna delle chiamate controlla il valore di ritorno.
    \item Riga 16: \texttt{length} non è dichiarata (\texttt{socklen\_t}); ed è un parametro valore-risultato, da reimpostare prima di ogni
          \texttt{accept}.
    \item Riga 19: il ciclo vuoto su \texttt{-1} gira a vuoto per sempre su un errore permanente; si ritenta solo su \texttt{EINTR}.
    \item Riga 22: il figlio legge dal socket \textbf{in ascolto} \texttt{ds\_sock} invece che da quello connesso \texttt{ds\_sock\_a};
          non controlla quanti byte ha letto né se il client ha chiuso (0); \texttt{buff} non è terminato.
    \item Riga 24: \texttt{strcmp} con \texttt{"quit"} fallisce se arriva \texttt{"quit\textbackslash n"}, o se \texttt{quit} arriva spezzato
          o insieme ad altri byte: una \texttt{read} non è un messaggio.
    \item Riga 25: scrive sul socket sbagliato, e 10 byte da una stringa di 5 (legge oltre la costante).
    \item Il figlio non chiude il socket in ascolto che ha ereditato.
\end{itemize}
\end{esercizio}

% ══════════════════════════════════════════════════════════════
\section{Processi}
% ══════════════════════════════════════════════════════════════

\begin{esercizio}{fork in un ciclo}
\begin{lstlisting}[language={[POSIX]C}]
int main(void) {
    for (int i = 0; i < 2; i++)
        if (fork() > 0) printf("parent %d\n", i);
        else            printf("child %d\n", i);
    sleep(10);
    return 0;
}
\end{lstlisting}
\begin{enumerate}[label=(\alph*)]
    \item Quante righe vengono stampate, e quanti processi esistono, contando il primo? Disegnare l'albero.
    \item Si può prevedere l'ordine delle righe?
    \item Un server esegue \texttt{if (fork() == 0) serve(c);} nel ciclo di \texttt{accept}, e \texttt{serve} ritorna quando il client se
          ne va. Che cosa succede poi in quel figlio, e quanti processi finiscono per attendere sullo stesso socket in ascolto?
\end{enumerate}

\svolgimento
\begin{enumerate}[label=(\alph*)]
    \item Con $i = 0$ due processi stampano una riga ciascuno; con $i = 1$ i due processi fanno ciascuno una \texttt{fork}, e quattro
          processi stampano. In tutto \textbf{6 righe} e \textbf{4 processi}: P genera C1 (con $i=0$) e C2 (con $i=1$); C1 genera C3 (con
          $i=1$).
    \item No: decide lo scheduler. Si sa solo che ogni processo stampa la sua riga per $i=0$ prima di quella per $i=1$.
    \item Il figlio esce da \texttt{serve}, torna in cima al ciclo e chiama \texttt{accept} come se fosse il padre: diventa un altro server.
          Dopo $k$ client ci sono fino a $k+1$ processi in attesa sullo stesso socket. Il figlio deve terminare con \texttt{\_exit} dopo
          \texttt{serve}.
\end{enumerate}
\end{esercizio}

\begin{esercizio}{Buffer e fork}
Che cosa stampa questo programma, e perché? Che cosa cambia se entrambi i \texttt{printf} diventano \texttt{write(1, ...)}?
\begin{lstlisting}[language={[POSIX]C}]
printf("A");
fork();
printf("B\n");
\end{lstlisting}

\svolgimento
Stampa \texttt{AB} due volte. \texttt{"A"} resta nel buffer di \texttt{stdio} (non c'è un a capo), la \texttt{fork} duplica anche il buffer,
e ciascun processo lo svuota con il proprio \texttt{"B\textbackslash n"}. Con \texttt{write} non c'è buffer in spazio utente:
\texttt{A} esce una volta sola prima della \texttt{fork}, e si vede \texttt{AB} seguito da \texttt{B}.
\end{esercizio}

\begin{esercizio}{Un file condiviso tra padre e figlio}
Scrivere un programma che apre un file, fa \texttt{fork}, e fa scrivere a padre e figlio cinque righe ciascuno.
\begin{enumerate}[label=(\alph*)]
    \item Come si intercalano le righe? Qualcuna ne sovrascrive un'altra?
    \item Spostare la \texttt{open} dopo la \texttt{fork}: che cosa cambia, e perché?
    \item Un server fa \texttt{fork} dopo \texttt{accept}: in quale dei due casi si trova il socket connesso?
\end{enumerate}

\svolgimento
\begin{enumerate}[label=(\alph*)]
    \item Con la \texttt{open} prima della \texttt{fork} i due descrittori puntano alla \textbf{stessa} descrizione di file aperto, con un
          solo offset condiviso: le righe si intercalano secondo lo scheduling, ma nessuna sovrascrive l'altra.
    \item Ognuno ha la propria descrizione e il proprio offset, entrambi a 0: si sovrascrivono a vicenda (a meno di \texttt{O\_APPEND}).
    \item Nel primo: padre e figlio condividono la stessa descrizione del socket connesso, ed è per questo che il padre deve chiuderlo.
\end{enumerate}
\end{esercizio}

% ══════════════════════════════════════════════════════════════
\section{Thread e sincronizzazione}
% ══════════════════════════════════════════════════════════════

\begin{esercizio}{Una variabile, due client}
Un server con thread passa ogni connessione al suo worker così:
\begin{lstlisting}[language={[POSIX]C}]
for (;;) {
    int c = accept(ls, NULL, NULL);
    pthread_create(&tid, NULL, serve, &c);
    pthread_detach(tid);
}

void *serve(void *arg) {
    int fd = *(int *) arg;
    /* ... parla con il client su fd, poi lo chiude ... */
}
\end{lstlisting}
\begin{enumerate}[label=(\alph*)]
    \item Due client si collegano a pochi microsecondi di distanza. Descrivere un ordine di eventi in cui entrambi i worker servono il
          secondo client. Che cosa succede al primo, e al suo descrittore?
    \item Correggerlo in due modi diversi.
\end{enumerate}

\svolgimento
\begin{enumerate}[label=(\alph*)]
    \item Il ciclo accetta il primo client ($c = 4$) e crea il worker 1, che non ha ancora letto \texttt{*arg}; accetta il secondo
          ($c = 5$) sulla \textbf{stessa variabile}; ora entrambi i worker leggono 5. Il primo client non viene mai servito e resta appeso;
          il descrittore 4 non viene mai chiuso (perdita), e il 5 viene chiuso due volte: la seconda \texttt{close} può chiudere un
          descrittore nel frattempo riassegnato a un altro client.
    \item Si passa a ogni thread una copia sua: \texttt{int *p = malloc(sizeof *p); *p = c;} e il worker fa \texttt{free}; oppure si passa
          il valore dentro il puntatore, \texttt{(void *)(intptr\_t) c}, e il worker lo riconverte. (Un terzo modo: un semaforo che il main
          attende finché il worker non ha copiato il valore.)
\end{enumerate}
\end{esercizio}

\begin{esercizio}{Una wait dentro un if}
Due worker eseguono questa funzione. Il main li crea, dorme un secondo, poi, tenendo \texttt{m}, pone \texttt{ready = 1} e chiama
\texttt{pthread\_cond\_broadcast(\&c)}.
\begin{lstlisting}[language={[POSIX]C}]
int ready = 0;

void *worker(void *arg) {
    int id = (int)(long) arg;
    pthread_mutex_lock(&m);
    if (!ready)
        pthread_cond_wait(&c, &m);
    printf("thread %d in the critical section\n", id);
    ready = 0;                  /* questo thread "consuma" la condizione */
    pthread_mutex_unlock(&m);
    return NULL;
}
\end{lstlisting}
(a) Che cosa stampa il programma? (b) Quando ciascun thread entra nella sezione critica, quanto vale \texttt{ready}? (c) Correggerlo.

\svolgimento
(a) Entrambi i thread stampano la loro riga: il broadcast li sveglia tutti e due. (b) Il primo che riprende il mutex trova
\texttt{ready = 1} e lo azzera; il secondo entra con \texttt{ready = 0}, perché con l'\texttt{if} non ricontrolla la condizione. (c)
\texttt{while (!ready) pthread\_cond\_wait(\&c, \&m);}: il secondo torna ad aspettare finché qualcuno non rimette \texttt{ready} a 1. Se si
vuole far passare un solo thread per segnale si usa \texttt{signal}; il \texttt{while} resta comunque necessario (risvegli spuri, e un
altro thread può aver consumato la condizione prima).
\end{esercizio}

\begin{esercizio}{Thread che si svegliano a turno}
Scrivere un programma C che crea tre thread, t1, t2 e t3, che devono eseguire in questo ordine: t1 stampa \texttt{t1 starts}, poi
\texttt{t1 ends}, e sveglia t2; t2 aspetta di essere svegliato, stampa \texttt{t2 starts} e \texttt{t2 ends}, e sveglia t3; t3 aspetta e
stampa \texttt{t3 starts} e \texttt{t3 ends}. Un solo intero globale per il turno, un mutex e una variabile di condizione; il main fa join
dei tre e distrugge ciò che ha creato. Prova: creare i thread nell'ordine t3, t2, t1 e aggiungere un \texttt{nanosleep} casuale prima di
ogni attesa: l'output non deve cambiare. Con una sola variabile di condizione, si deve usare \texttt{signal} o \texttt{broadcast}?

\svolgimento
\begin{lstlisting}[language={[POSIX]C}]
static int turn = 1;
static pthread_mutex_t m = PTHREAD_MUTEX_INITIALIZER;
static pthread_cond_t  c = PTHREAD_COND_INITIALIZER;

static void *worker(void *arg) {
    int me = (int)(long) arg;
    unsigned seed = (unsigned) time(NULL) ^ (unsigned) me;
    struct timespec t = { 0, (rand_r(&seed) % 200) * 1000000L };
    nanosleep(&t, NULL);                       /* ritardo casuale */
    pthread_mutex_lock(&m);
    while (turn != me)                         /* while, non if */
        pthread_cond_wait(&c, &m);
    printf("t%d starts\nt%d ends\n", me, me);
    turn++;
    pthread_cond_broadcast(&c);                /* sveglia tutti: decide turn */
    pthread_mutex_unlock(&m);
    return NULL;
}

int main(void) {
    pthread_t t[3];
    for (int i = 3; i >= 1; i--)               /* creati come t3, t2, t1 */
        pthread_create(&t[i - 1], NULL, worker, (void *)(long) i);
    for (int i = 0; i < 3; i++)
        pthread_join(t[i], NULL);
    pthread_mutex_destroy(&m);
    pthread_cond_destroy(&c);
    return 0;
}
\end{lstlisting}
Serve \texttt{broadcast}. Con una sola variabile di condizione \texttt{signal} sveglia un thread qualsiasi tra quelli in attesa: se t1
sveglia t3 invece di t2, t3 trova che non è il suo turno e torna ad aspettare, t2 non viene mai svegliato e il programma si blocca. Con
\texttt{broadcast} si svegliano tutti, e il \texttt{while} lascia passare solo quello giusto.
\end{esercizio}

\begin{esercizio}{Un client che invia e ascolta}
Un client UDP ha un solo socket verso un server di eco su \texttt{127.0.0.1:5000}. Avvia \texttt{sender}, poi \texttt{receiver}, e fa join
di entrambi.
\begin{lstlisting}[language={[POSIX]C}]
void *sender(void *arg) {
    pthread_mutex_lock(&mtx);
    sendto(fd, "A", 1, 0, (struct sockaddr *) &srv, sizeof srv);
    sendto(fd, "B", 1, 0, (struct sockaddr *) &srv, sizeof srv);
    pthread_mutex_unlock(&mtx);
    return NULL;
}
void *receiver(void *arg) {
    char buf[64];
    pthread_mutex_lock(&mtx);
    ssize_t n = recvfrom(fd, buf, sizeof buf - 1, 0, NULL, NULL);
    pthread_mutex_unlock(&mtx);
    if (n > 0) { buf[n] = '\0'; printf("got %s\n", buf); }
    return NULL;
}
\end{lstlisting}
\begin{enumerate}[label=(\alph*)]
    \item Trovare un ordine di eventi in cui il programma non termina mai. Quale thread tiene che cosa, e che cosa aspetta?
    \item Che cosa succede se il server di eco non è in esecuzione?
    \item Che cosa protegge qui il mutex? Qualcosa in queste due funzioni ne ha bisogno mentre l'I/O è in corso?
    \item Correggere il progetto in modo che il client termini sempre, qualunque thread parta per primo.
\end{enumerate}

\svolgimento
\begin{enumerate}[label=(\alph*)]
    \item Il receiver parte per primo, prende il mutex e si blocca in \texttt{recvfrom}; il sender si blocca sul mutex e non invierà mai
          nulla. Il receiver tiene il mutex e aspetta un datagramma che può arrivare solo dopo che il sender avrà avuto il mutex: stallo.
    \item Anche nell'ordine buono, i datagrammi vengono scartati (l'ICMP \emph{port unreachable} non viene riferito a un socket non
          connesso) e il receiver resta in \texttt{recvfrom} per sempre.
    \item Niente: le chiamate sullo stesso socket da thread diversi sono già serializzate dal kernel, e non c'è stato condiviso in spazio
          utente. Tenere un mutex durante un'operazione bloccante è sbagliato in ogni caso.
    \item Si toglie il mutex attorno all'I/O; il receiver imposta \texttt{SO\_RCVTIMEO} e legge al più due risposte (o fino al timeout);
          meglio ancora si fa \texttt{connect} sul socket UDP, così un server assente produce \texttt{ECONNREFUSED} invece di un'attesa
          infinita.
\end{enumerate}
\end{esercizio}

% ══════════════════════════════════════════════════════════════
\section{Server concorrenti e multiplexing}
% ══════════════════════════════════════════════════════════════

\begin{esercizio}{Il padre ha dimenticato una close}
In questo server di eco con \texttt{fork}, \texttt{serve} fa l'eco delle righe finché il client non invia \texttt{quit}, poi ritorna.
\begin{lstlisting}[language={[POSIX]C}]
for (;;) {
    int c = accept(ls, NULL, NULL);
    if (fork() == 0) {
        close(ls);
        serve(c);
        close(c);
        _exit(0);
    }
    /* qui il padre dovrebbe fare close(c), e non lo fa */
}
\end{lstlisting}
\begin{enumerate}[label=(\alph*)]
    \item Un client invia \texttt{quit}. Il figlio chiude il suo descrittore ed esce. La \texttt{read} successiva del client restituisce 0?
          Perché?
    \item Che cosa mostra \texttt{ss -tn} sul server dopo?
    \item Dopo qualche migliaio di client, che cos'altro si esaurisce?
\end{enumerate}

\svolgimento
\begin{enumerate}[label=(\alph*)]
    \item No: il padre ha ancora un descrittore aperto sullo stesso socket. Un socket viene chiuso, e il FIN parte, solo quando si chiude
          l'ultimo descrittore che vi fa riferimento. Il client resta bloccato.
    \item La connessione resta \texttt{ESTAB}; se poi il client chiude dal suo lato, sul server passa in \texttt{CLOSE-WAIT} e ci resta per
          sempre. Si accumula una riga per ogni client servito.
    \item I \textbf{descrittori} del padre (il limite tipico è 1024): \texttt{accept} comincia a fallire con \texttt{EMFILE}, e il server
          smette di accettare.
\end{enumerate}
\end{esercizio}

\begin{esercizio}{Trovare gli errori in un ciclo select}
Un client deve inoltrare al server ciò che l'utente scrive, e stampare ciò che il server invia, qualunque arrivi prima.
\begin{lstlisting}[language={[POSIX]C}]
fd_set rset;
FD_ZERO(&rset);
FD_SET(0, &rset);
FD_SET(sock, &rset);
for (;;) {
    select(sock, &rset, NULL, NULL, NULL);
    if (FD_ISSET(0, &rset))    { /* legge una riga da 0, la invia */ }
    if (FD_ISSET(sock, &rset)) { /* legge da sock, la stampa */ }
}
\end{lstlisting}
Gli utenti segnalano due sintomi: niente di ciò che arriva dal server viene mai stampato, e dopo la prima riga il programma si comporta
in modo strano. Trovare i due errori che li causano, e una terza cosa che il ciclo non controlla mai.

\svolgimento
\begin{itemize}
    \item Il primo argomento di \texttt{select} è il \textbf{massimo descrittore più uno}: con \texttt{sock} invece di \texttt{sock + 1} il
          socket non viene mai controllato, e il server non viene mai stampato.
    \item \texttt{select} \textbf{modifica} l'insieme, lasciandovi solo i descrittori pronti: va ricostruito (\texttt{FD\_ZERO} e
          \texttt{FD\_SET}) a ogni giro, dentro il ciclo. Altrimenti dopo il primo evento si ascolta solo il descrittore che era pronto.
    \item Non si controlla mai il valore di ritorno di \texttt{select} (\texttt{-1}, \texttt{EINTR}) né la fine dei flussi: una
          \texttt{read} che restituisce 0 su \texttt{sock} (il server ha chiuso) o su 0 (fine dell'input) fa girare il ciclo a vuoto per
          sempre.
\end{itemize}
\end{esercizio}

\begin{esercizio}{Il server che la lezione dava per scontato}
Scrivere il server dietro il client multiplexato della lezione: serve un client alla volta, in un solo thread; fa l'eco di ogni riga
ricevuta; ogni 10 secondi invia anche una riga \texttt{x}, qualunque cosa stia facendo il client; prende la porta dalla riga di comando
(\texttt{./tick\_server 6000}). Usare \texttt{select} con un timeout: i dieci secondi devono restare dieci anche quando le righe continuano
ad arrivare. Prova: collegarsi con \texttt{nc 127.0.0.1 6000}, non scrivere nulla e contare le \texttt{x} in un minuto; poi scrivere una
riga ogni due secondi e contare di nuovo.

\svolgimento
Il punto è che il timeout di \texttt{select} non è sempre 10 secondi, ma il tempo che \textbf{manca} alla prossima scadenza, ricalcolato a
ogni giro:
\begin{lstlisting}[language={[POSIX]C}]
static double now(void) {
    struct timespec t; clock_gettime(CLOCK_MONOTONIC, &t);
    return t.tv_sec + t.tv_nsec / 1e9;
}

for (;;) {                                         /* un client alla volta */
    int c = accept(ls, NULL, NULL);
    double next = now() + 10;                      /* prossima x */
    char buf[1024];
    for (;;) {
        fd_set r; FD_ZERO(&r); FD_SET(c, &r);
        double left = next - now();
        if (left < 0) left = 0;
        struct timeval tv = { (time_t) left, (suseconds_t) ((left - (time_t) left) * 1e6) };
        int k = select(c + 1, &r, NULL, NULL, &tv);
        if (k < 0) { if (errno == EINTR) continue; break; }
        if (now() >= next) {                       /* scadenza raggiunta */
            write(c, "x\n", 2);
            next += 10;                            /* non now() + 10: niente deriva */
        }
        if (k > 0 && FD_ISSET(c, &r)) {
            ssize_t n = read(c, buf, sizeof buf);
            if (n <= 0) break;                     /* il client se ne è andato */
            write(c, buf, (size_t) n);             /* eco */
        }
    }
    close(c);
}
\end{lstlisting}
In entrambe le prove si contano \textbf{6} \texttt{x} al minuto. Con un timeout fisso di 10 secondi, ogni riga in arrivo farebbe ripartire
l'attesa, e scrivendo ogni due secondi non si vedrebbe mai una \texttt{x}.
\end{esercizio}

\begin{esercizio}{Un figlio che ascolta due sorgenti}
Scrivere un programma C che crea una pipe e fa \texttt{fork}. Il padre invia al figlio dieci stringhe attraverso la pipe, a intervalli
casuali tra uno e tre secondi, poi chiude il proprio estremo e aspetta il figlio. Il figlio aspetta sia sulla pipe sia sulla tastiera con
\texttt{select}, stampa \texttt{[PARENT] <testo>} per ciò che arriva dalla pipe e \texttt{[USER] <testo>} per ciò che scrive l'utente, e
termina quando il padre ha chiuso la pipe. Come fa il figlio a sapere che il padre ha chiuso la pipe, e che cosa succederebbe se il figlio
dimenticasse di chiudere la propria copia dell'estremo di scrittura?

\svolgimento
\begin{lstlisting}[language={[POSIX]C}]
int p[2];
pipe(p);
pid_t pid = fork();
if (pid == 0) {                                    /* figlio */
    close(p[1]);                                   /* essenziale, vedi sotto */
    char buf[256];
    int kbd = 1;                                   /* la tastiera è ancora aperta */
    for (;;) {
        fd_set r; FD_ZERO(&r); FD_SET(p[0], &r);
        if (kbd) FD_SET(0, &r);
        if (select(p[0] + 1, &r, NULL, NULL, NULL) < 0) continue;
        if (FD_ISSET(p[0], &r)) {
            ssize_t n = read(p[0], buf, sizeof buf - 1);
            if (n <= 0) break;                     /* EOF: il padre ha chiuso */
            buf[n] = '\0'; printf("[PARENT] %s", buf);
        }
        if (kbd && FD_ISSET(0, &r)) {
            ssize_t n = read(0, buf, sizeof buf - 1);
            if (n <= 0) kbd = 0;                   /* Ctrl-D: smette di ascoltarla */
            else { buf[n] = '\0'; printf("[USER] %s", buf); }
        }
        fflush(stdout);
    }
    _exit(0);
}
close(p[0]);                                       /* padre */
srand((unsigned) getpid());
for (int i = 1; i <= 10; i++) {
    sleep(1 + rand() % 3);
    char s[32];
    int len = snprintf(s, sizeof s, "stringa %d\n", i);
    write(p[1], s, (size_t) len);
}
close(p[1]);
waitpid(pid, NULL, 0);
\end{lstlisting}
Il figlio lo sa quando \texttt{read} sulla pipe restituisce \textbf{0}, cioè quando non esiste più alcun descrittore aperto
sull'estremo di scrittura. Se il figlio non chiudesse la propria copia, un estremo di scrittura resterebbe aperto (il suo), l'EOF non
arriverebbe mai, il figlio aspetterebbe per sempre e il padre resterebbe bloccato in \texttt{waitpid}.
\end{esercizio}

\begin{esercizio}{Dove aspetta il settimo client?}
Il ciclo di \texttt{accept} di un server mette ogni nuova connessione in una coda limitata di 4 posti; 2 thread worker prendono le
connessioni dalla coda e le servono. Quando la coda è piena, il ciclo di \texttt{accept} si blocca finché non si libera un posto. Servire
un client richiede 10 secondi. Il socket in ascolto ha un backlog di 5: il kernel tiene al più cinque connessioni completate che nessuno ha
ancora accettato.
\begin{enumerate}[label=(\alph*)]
    \item Dieci client si collegano entro un secondo. Un secondo dopo, dove si trova ciascuno: servito, in coda, nelle mani del ciclo di
          \texttt{accept}, nella coda del kernel, o rifiutato?
    \item La \texttt{connect} del decimo client è riuscita? Che cosa vede se invia una richiesta ora?
    \item Che cosa cambia se, invece di bloccarsi con la coda piena, il ciclo di \texttt{accept} chiude subito la nuova connessione?
\end{enumerate}

\svolgimento
\begin{enumerate}[label=(\alph*)]
    \item Client 1 e 2: serviti dai due worker. Client 3--6: nella coda di 4 posti. Client 7: accettato, nelle mani del ciclo di
          \texttt{accept}, bloccato nel tentativo di metterlo in coda. Client 8, 9 e 10: nella coda del kernel (3 dei 5 posti). Nessuno
          rifiutato.
    \item Sì: l'handshake lo completa il kernel, senza aspettare \texttt{accept}. La richiesta viene messa nel buffer di ricezione del
          socket e il client non riceve nulla finché un worker non lo prende: decine di secondi (dopo 10 s si libera il primo worker, e
          così via).
    \item I client 7--10 vengono accettati e chiusi subito: vedono la connessione chiusa (la loro \texttt{read} restituisce 0, o ricevono
          un RST se avevano inviato dati) invece di attendere. Un fallimento rapido e visibile al posto di un'attesa lunga.
\end{enumerate}
\end{esercizio}

% ══════════════════════════════════════════════════════════════
\section{UDP, broadcast e multicast}
% ══════════════════════════════════════════════════════════════

\begin{esercizio}{Chi lo riceve?}
L'host A, \texttt{192.168.1.23/24}, invia datagrammi UDP alla porta 6000. Sulla stessa LAN, B (\texttt{192.168.1.40}) ha un socket legato
alla porta 6000; C (\texttt{192.168.1.50}) ne ha uno legato alla 6000 e iscritto al gruppo \texttt{239.1.2.3}. D, su
\texttt{192.168.2.0/24} dietro un router, è legato alla 6000 e iscritto allo stesso gruppo.
\begin{enumerate}[label=(\alph*)]
    \item Per ogni indirizzo di destinazione, chi riceve il datagramma: \texttt{255.255.255.255}; \texttt{192.168.1.255};
          \texttt{239.1.2.3} con TTL 1; \texttt{239.1.2.3} con TTL 8, attraverso un router che inoltra il multicast.
    \item Che cosa deve impostare A sul socket prima di inviare a \texttt{255.255.255.255}, e che cosa succede se non lo fa?
    \item A riceve i propri datagrammi in qualcuno dei casi?
\end{enumerate}

\svolgimento
\begin{enumerate}[label=(\alph*)]
    \item \texttt{255.255.255.255}: B e C (il broadcast limitato non attraversa il router). \texttt{192.168.1.255}: B e C (broadcast
          diretto della propria subnet). \texttt{239.1.2.3} con TTL 1: solo C, l'unico iscritto sulla LAN; il router non inoltra un TTL 1.
          Con TTL 8: C e D.
    \item L'opzione \texttt{SO\_BROADCAST}; senza, \texttt{sendto} fallisce con \texttt{EACCES}.
    \item Solo se A ha a sua volta un socket legato alla porta 6000: in quel caso riceve il proprio broadcast, e il proprio multicast se è
          iscritto al gruppo (\texttt{IP\_MULTICAST\_LOOP} è attivo per default). Il socket da cui invia, senza quella porta, non riceve
          nulla.
\end{enumerate}
\end{esercizio}

\begin{esercizio}{Confini dei messaggi}
I socket TCP trattano i dati come un flusso di byte; quelli UDP conservano i confini dei messaggi. Dare un vantaggio e uno svantaggio di
un'API che conserva i confini dei messaggi applicativi.

\svolgimento
Vantaggio: l'applicazione non deve delimitare i messaggi; una \texttt{recvfrom} restituisce esattamente un messaggio. Svantaggio: ogni
messaggio deve stare in un datagramma (limite di dimensione, e frammentazione IP oltre la MTU), e un buffer troppo piccolo tronca il
messaggio invece di lasciarne il resto per la lettura successiva; non si può trasmettere a flusso un dato più grande.
\end{esercizio}

\begin{esercizio}{Un protocollo sopra TCP}
\begin{enumerate}[label=(\alph*)]
    \item Un client chiama \texttt{write(fd, "HELLO\textbackslash n", 6)} e poi \texttt{write(fd, "WORLD\textbackslash n", 6)}; il server
          chiama una volta \texttt{read(fd, buf, 100)}. Elencare tutti i risultati che la \texttt{read} può restituire.
    \item Ogni messaggio è delimitato da una lunghezza di 4 byte big-endian seguita dal corpo. Il server legge, in esadecimale,
          \texttt{00 00 00 03 61 62 63 00 00 00 00 00 00 00 02 68 69}. Dividerlo in messaggi.
    \item Che cosa deve fare il server se legge \texttt{7F FF FF FF} come lunghezza, e prima di quale chiamata?
\end{enumerate}

\svolgimento
\begin{enumerate}[label=(\alph*)]
    \item Qualsiasi prefisso non vuoto dei 12 byte \texttt{HELLO\textbackslash nWORLD\textbackslash n}: da 1 a 12 byte. In pratica di solito
          6 o 12, ma nessuno lo garantisce.
    \item Tre messaggi: \texttt{"abc"} (lunghezza 3), un messaggio \textbf{vuoto} (lunghezza 0), \texttt{"hi"} (lunghezza 2).
    \item $\texttt{0x7FFFFFFF} \approx 2$ GiB: oltre il massimo che il protocollo ammette. Il server deve rifiutarlo e chiudere la connessione
          \textbf{prima} di allocare memoria (\texttt{malloc}) e di leggere il corpo: altrimenti un solo client può esaurirgli la memoria.
\end{enumerate}
\end{esercizio}

% ══════════════════════════════════════════════════════════════
\section{Programmi da scrivere e da osservare}
% ══════════════════════════════════════════════════════════════

\begin{esercizio}{Il vostro server come proxy del browser}
Scrivere un server TCP in C che accetta una connessione alla volta e stampa sullo standard output ogni riga che riceve, con il numero di
riga (\texttt{./dump 8080}). Impostare il proxy HTTP di un browser sul proprio host e sulla porta 8080, e aprire una pagina servita in
\texttt{http://}.
\begin{enumerate}[label=(\alph*)]
    \item In che cosa differisce la prima riga di ogni richiesta dal GET scritto a mano nel primo pacchetto di esercizi?
    \item Ricaricare la pagina. Il browser invia richieste condizionali (\texttt{If-Modified-Since:}, \texttt{If-None-Match:})?
    \item Il browser continua ad aspettare. Che cosa dovrebbe fare il server per soddisfarlo?
\end{enumerate}
Le pagine servite in \texttt{https://} compaiono come un'unica riga \texttt{CONNECT}: è tutto ciò che un proxy può vederne.

\svolgimento
\begin{lstlisting}[language={[POSIX]C}]
int ls = socket(AF_INET, SOCK_STREAM, 0);
int one = 1;
setsockopt(ls, SOL_SOCKET, SO_REUSEADDR, &one, sizeof one);
struct sockaddr_in a = { .sin_family = AF_INET, .sin_port = htons((uint16_t) atoi(argv[1])),
                         .sin_addr.s_addr = htonl(INADDR_ANY) };
bind(ls, (struct sockaddr *) &a, sizeof a);
listen(ls, 5);
for (;;) {
    int c = accept(ls, NULL, NULL);
    FILE *in = fdopen(c, "r");                     /* fgets legge per righe */
    char line[4096];
    int k = 0;
    while (fgets(line, sizeof line, in))
        printf("%4d %s", ++k, line);
    fclose(in);                                    /* chiude anche c */
}
\end{lstlisting}
\begin{enumerate}[label=(\alph*)]
    \item Verso un proxy la riga di richiesta contiene l'\textbf{URL assoluto}, \texttt{GET http://host/percorso HTTP/1.1}, non solo il
          percorso: il proxy deve sapere a quale server rivolgersi.
    \item Sì, se il browser ha la pagina in cache: \texttt{If-Modified-Since:} e/o \texttt{If-None-Match:} con l'ETag.
    \item Rispondere con un messaggio HTTP: cioè fare davvero il proxy, inoltrando la richiesta al server d'origine e restituendo la
          risposta (o almeno inviare una risposta, per esempio un \texttt{502}).
\end{enumerate}
\end{esercizio}

\begin{esercizio}{Un pinger UDP}
Scrivere due programmi: \texttt{./udp\_echo 7000} rimanda ogni datagramma che riceve, ma ne scarta ciascuno con probabilità 0,3;
\texttt{./udp\_ping <host> 7000} invia dieci datagrammi \texttt{ping <n> <tempo>}, uno al secondo, aspetta al più un secondo per ogni
risposta (\texttt{SO\_RCVTIMEO}) e stampa l'RTT in millisecondi, oppure \texttt{lost}; alla fine minimo, media e massimo dell'RTT e il
tasso di perdita. Prova su \texttt{127.0.0.1}, poi tra due macchine, con una cattura in corso.
\begin{enumerate}[label=(\alph*)]
    \item C'è un datagramma sul filo per ogni \texttt{sendto}?
    \item Quando una risposta si perde, viaggia qualcosa a dirlo?
    \item Una risposta che arriva dopo il proprio timeout: che cosa ne fa il client, e come se ne accorge?
\end{enumerate}

\svolgimento
\begin{lstlisting}[language={[POSIX]C}]
/* udp_echo: il ciclo principale */
for (;;) {
    char buf[1500];
    struct sockaddr_in from;
    socklen_t fl = sizeof from;
    ssize_t n = recvfrom(s, buf, sizeof buf, 0, (struct sockaddr *) &from, &fl);
    if (n < 0 || rand() % 10 < 3) continue;        /* ne perde il 30% */
    sendto(s, buf, (size_t) n, 0, (struct sockaddr *) &from, fl);
}

/* udp_ping: un giro per ogni n */
struct timeval to = { 1, 0 };
setsockopt(s, SOL_SOCKET, SO_RCVTIMEO, &to, sizeof to);
for (int i = 1; i <= 10; i++) {
    char msg[64], buf[64];
    double t0 = now_ms();                          /* clock_gettime(CLOCK_MONOTONIC) */
    int len = snprintf(msg, sizeof msg, "ping %d %.3f", i, t0);
    sendto(s, msg, (size_t) len, 0, (struct sockaddr *) &srv, sizeof srv);
    for (;;) {
        ssize_t n = recvfrom(s, buf, sizeof buf - 1, 0, NULL, NULL);
        if (n < 0) { printf("%d: lost\n", i); lost++; break; }   /* timeout */
        buf[n] = '\0';
        int seq;
        if (sscanf(buf, "ping %d", &seq) == 1 && seq == i) {
            double rtt = now_ms() - t0;
            printf("%d: %.3f ms\n", i, rtt);
            /* ... aggiorna minimo, massimo e somma ... */
            break;
        }                                          /* risposta vecchia: la si scarta */
    }
    sleep(1);
}
\end{lstlisting}
\begin{enumerate}[label=(\alph*)]
    \item Sì: un datagramma per ogni \texttt{sendto} (i messaggi sono piccoli, nessuna frammentazione).
    \item Niente: UDP non ha riscontri. La perdita si scopre solo perché scade il timeout del client.
    \item Arriva mentre il client aspetta la risposta a un ping successivo: lo riconosce dal numero di sequenza nel payload, diverso da
          quello atteso, e la scarta (o la segnala come tardiva). Senza numero di sequenza la scambierebbe per la risposta giusta, con un RTT
          sbagliato.
\end{enumerate}
\end{esercizio}

\begin{esercizio}{Leggere le proprie connessioni}
Con una pipeline di comandi, elencare tutte le connessioni TCP in stato \texttt{ESTAB} e stampare solo l'indirizzo IP remoto, senza
porta: da \texttt{192.168.1.10:45678} deve stampare \texttt{192.168.1.10}. Suggerimento: \texttt{ss -tn}, e \texttt{cut} con \texttt{:}
come delimitatore.
\begin{enumerate}[label=(\alph*)]
    \item Provarla con un browser aperto su qualche sito.
    \item Perché il semplice \texttt{cut} si rompe sugli indirizzi IPv6, come \texttt{[2a00:1450:4002::65]:443}? Cambiare la pipeline in
          modo che gestisca entrambi.
\end{enumerate}

\svolgimento
\begin{lstlisting}[language=bash]
$ ss -tn | grep ESTAB | awk '{print $5}' | cut -d: -f1
$ ss -tn | grep ESTAB | awk '{print $5}' | sed -E 's/:[0-9]+$//; s/^\[//; s/\]$//'
\end{lstlisting}
La quinta colonna è l'indirizzo remoto. Un indirizzo IPv6 contiene a sua volta i due punti, quindi \texttt{cut -d: -f1} si ferma al primo
gruppo (e lascia la parentesi quadra). La seconda versione toglie solo l'\textbf{ultimo} \texttt{:porta}, poi le parentesi quadre, e
funziona per entrambe le famiglie.
\end{esercizio}

\begin{esercizio}{I vostri programmi sul filo}
Un server e un client TCP vostri sulla porta 5000, con una cattura sull'interfaccia di loopback e filtro \texttt{tcp.port == 5000}.
\begin{enumerate}[label=(\alph*)]
    \item Quale chiamata fa partire il primo segmento? Dove si trova il server quando i primi tre segmenti sono passati: già dentro
          \texttt{accept}, o non ancora?
    \item Il client scrive 10 byte in una chiamata. Quanti segmenti li portano, e che cosa torna indietro?
    \item Ora ne scrive 100\,000 in una chiamata. Quanti segmenti, qual è il payload più grande, e quanti segmenti tornano?
    \item Subito dopo l'uscita del client, eseguire \texttt{ss -tan}: quali nomi di stato compaiono per la porta 5000, e da quale lato?
\end{enumerate}

\svolgimento
\begin{enumerate}[label=(\alph*)]
    \item \texttt{connect} invia il SYN. L'handshake lo completano i kernel: il server può essere dentro \texttt{accept} o non esserci
          ancora arrivato, e la connessione aspetta nella coda di backlog.
    \item Un segmento con 10 byte di payload; torna un ACK (puro, o portato da una risposta).
    \item Su loopback la MTU è 65\,536 byte, quindi l'MSS è circa 65\,483: bastano 2 segmenti (il più grande vicino a 65\,483 byte), e
          tornano uno o due ACK. Su Ethernet, con MSS 1448, sarebbero circa 70 segmenti, con un ACK ogni due.
    \item Chi chiude per primo (il client) resta in \texttt{TIME-WAIT} per un minuto circa; il server, se ha già chiuso il suo lato, non ha
          più righe (o mostra brevemente \texttt{LAST-ACK}); il socket in ascolto resta \texttt{LISTEN}.
\end{enumerate}
\end{esercizio}

\begin{esercizio}{Datagrammi, sul filo}
Con un mittente e un ricevente UDP vostri, e una cattura in corso:
\begin{enumerate}[label=(\alph*)]
    \item Una \texttt{sendto} di 3000 byte, poi tre da 1000. Quanti datagrammi si vedono su loopback? Tra due macchine della stessa LAN se
          ne vedrà qualcosa in più: annotare che cosa.
    \item Il ricevente chiama \texttt{recvfrom} con un buffer di 100 byte su un datagramma di 1000. Che cosa restituisce, e dove sono gli
          altri 900 byte dopo la chiamata?
    \item Il mittente invia 2000 datagrammi più in fretta che può; il ricevente ne legge uno al millisecondo. Quanti ne riceve, e dove si
          perdono gli altri? TCP li perderebbe?
\end{enumerate}

\svolgimento
\begin{enumerate}[label=(\alph*)]
    \item Su loopback quattro datagrammi. Tra due macchine il datagramma da 3000 byte supera la MTU di 1500 e viene \textbf{frammentato} da
          IP in tre frammenti: è la frammentazione del livello di rete.
    \item Restituisce 100: i primi 100 byte. Gli altri 900 sono \textbf{scartati} (il flag \texttt{MSG\_TRUNC} lo segnalerebbe): un
          datagramma si legge tutto in una volta o mai.
    \item Molti meno di 2000: quanti ne stanno nel buffer di ricezione del socket più quelli letti nel frattempo. Gli altri si perdono nel
          kernel del ricevente, a buffer pieno, non sul filo. TCP non li perderebbe: il controllo di flusso rallenterebbe il mittente.
\end{enumerate}
\end{esercizio}
